\section{Пробелы отчётности в исходной оценке}
\label{sec:gaps}

Собрано во время полного извлечения текста исходной статьи, каждый пункт
можно проверить по ней напрямую\origin. Они записаны здесь потому, что
несколько из них — причина, по которой воспроизведению приходится делать
собственный выбор, а читатель должен знать, какой из этих выборов был
вынужденным.

\begin{enumerate}
\item \textbf{Схемы.} Дословный список полей дан только для одного из
      четырёх доменов оценки. Для двух других схему можно вывести лишь из
      единственной примерной траектории, а для четвёртого не дано вообще
      ничего. Метод проектирования схемы нигде не описан, анализа
      чувствительности тоже нет. Воспроизводящему поэтому приходится
      изобретать схему самому, что само по себе источник расхождения — и
      мотивация трактовать проектирование схемы как фактор, а не как
      константу.

\item \textbf{Кэширование и деньги.} Полнотекстовый поиск не находит ни
      одного вхождения слов cache, billing или иной денежной терминологии.
      Все экономические утверждения даны в логических числах токенов.
      Статья не кратка на этот счёт — она вообще молчит по этому вопросу.

\item \textbf{Дисперсия и повторения.} Таблицы устойчивости к шуму и
      восстановления состояния приводят точечные значения без указания
      стандартного отклонения или числа повторений.

\item \textbf{Контроль сжатия недоопределён.} Базовую линию обученного
      сжатия приводят без коэффициента сжатия и без указания варианта.
      Бюджет усечения в одном месте указан в символах, а в другом — в
      токенах.

\item \textbf{Статистическое утверждение ничем не показано.} Значимость
      заявлена одним общим парным $t$-тестом при $p<0.01$ для горизонтов
      $\ge 50$ — без значений по отдельным сравнениям и без поправки на
      множественные сравнения.

\item \textbf{Незамеченный нуль.} В одной конфигурации с открытыми весами
      на самом длинном из приводимых горизонтов рантайм со структурированным
      состоянием и стабильная базовая линия набирают идентично. Таблица это
      число содержит; текст его не обсуждает.

\item \textbf{Авторитет над числом шагов.} Горизонт — это подсчёт внешне
      генерируемых событий, и ни одно условие не варьирует, кто управляет
      числом шагов. Это свойство разобрано в \S\ref{sec:repro}, и оно —
      самое значимое из всех различий между той оценкой и нашей.
\end{enumerate}
